% --- Archon physics formalization source begin ---
% archon:physics
% archon:covers PhyXMiniProblems/problem_phyx_mini_0785.lean
% archon:source-report reports/phyx_mini/problem_phyx_mini_0785.source.json

\chapter{Physics problem phyx\_mini\_0785}
\label{ch:PhyXMiniProblems_problem_phyx_mini_0785}

\paragraph{Problem source.}
\#\# Physical scenario

A block with mass \textbackslash{}(m = 5.00\textbackslash{}text\{ kg\}\textbackslash{}) slides down a surface inclined \textbackslash{}(36.9\textasciicircum{}\{\textbackslash{}circ\}\textbackslash{}) to the horizontal as shown in figure. The coefficient of kinetic friction is \textbackslash{}(0.25\textbackslash{}). A string attached to the block is wrapped around a flywheel on a fixed axis at \textbackslash{}(O\textbackslash{}). The flywheel has mass \textbackslash{}(25.0\textbackslash{}text\{ kg\}\textbackslash{}) and moment of inertia \textbackslash{}(0.500\textbackslash{}text\{ kg\}\textbackslash{}cdot\textbackslash{}text\{m\}\textasciicircum{}2\textbackslash{}) with respect to the axis of rotation. The string pulls without slipping at a perpendicular distance of \textbackslash{}(0.200\textbackslash{}text\{ m\}\textbackslash{}) from that axis.

\#\# Question

What is the tension in the string?

\#\# Answer choices

- A: A: 10.8N

- B: B: 14.0N

- C: C: 15.8N

- D: D: 21.0N

\#\# Figure caption (auxiliary; use the image as primary evidence)

The image depicts a physics diagram involving an inclined plane and a pulley system. The key elements in the image include:

- A triangular inclined plane with a slope labeled at an angle of 36.9° relative to the horizontal.
- An orange block resting on the inclined plane, labeled with a mass of 5.00 kg.
- A pulley located at the top of the inclined plane, with point "O" indicated alongside it.
- A rope or string connected from the pulley to the block, suggesting a tension or force scenario.

The setup likely represents a classical mechanics problem related to forces and motion on an inclined plane.

\#\# Dataset metadata

Dynamics; Physical Model Grounding Reasoning, Multi-Formula Reasoning

\paragraph{Recorded answer/context.}
B: B: 14.0N

\paragraph{Figure/image path.}
/root/proposal\_for\_physic/science-mango/phyx\_mini\_run/phyx\_data/test\_image/785.png

\paragraph{Formalization target.}
create a compiling Lean file with sorry bodies at `PhyXMiniProblems/problem\_phyx\_mini\_0785.lean`. The Lean declarations must preserve the physical quantities, dimensions or dimensional roles, figure labels, governing-law hypotheses, and final relation expressed by this problem.
Use Mathlib/Physlib names found through LeanExplore where available. If a physics API is missing, introduce faithful local abstractions rather than scalar placeholder aliases.

\begin{theorem}[Physics formalization target]
\label{thm:physics:phyx_mini_0785:target}
\lean{PhyXMiniProblems.ProblemPhyXMini0785.problem_phyx_mini_0785}
\uses{def:physics:phyx-mini-0785:phyxminiproblems-problemphyxmini0785-forceinnewtons, def:physics:phyx-mini-0785:phyxminiproblems-problemphyxmini0785-inclineflywheelsetup, def:physics:phyx-mini-0785:phyxminiproblems-problemphyxmini0785-matchesproblemreadouts, def:physics:phyx-mini-0785:phyxminiproblems-problemphyxmini0785-matchessuppliedfigure, def:physics:phyx-mini-0785:phyxminiproblems-problemphyxmini0785-usestextbookgravityandanglecalibration, def:physics:phyx-mini-0785:phyxminiproblems-problemphyxmini0785-hasphysicalparameters, def:physics:phyx-mini-0785:phyxminiproblems-problemphyxmini0785-satisfiesinclineflywheeldynamics, def:physics:phyx-mini-0785:phyxminiproblems-problemphyxmini0785-displayedtensioninnewtons, def:physics:phyx-mini-0785:phyxminiproblems-problemphyxmini0785-recordeddatasetanswer, def:physics:phyx-mini-0785:phyxminiproblems-problemphyxmini0785-isuniqueclosestdisplayedtension}
The assigned autoformalize agent should translate this physics problem into Lean declarations in the covered file, with theorem and lemma proof bodies written as `by sorry`.
\end{theorem}
\begin{proof}
This is an autoformalization task, not a proof task. Produce statements that can later be proved by the physics prover without weakening the physical model.
\end{proof}

% --- Archon declaration topology begin ---
\section{Declaration topology}
\begin{definition}[Mass Quantity]
  \label{def:physics:phyx-mini-0785:phyxminiproblems-problemphyxmini0785-massquantity}
  \lean{PhyXMiniProblems.ProblemPhyXMini0785.MassQuantity}
  A nonnegative physical mass
\end{definition}
\begin{proof}
  This is the declared model definition of Mass Quantity.
\end{proof}
\begin{definition}[Length Quantity]
  \label{def:physics:phyx-mini-0785:phyxminiproblems-problemphyxmini0785-lengthquantity}
  \lean{PhyXMiniProblems.ProblemPhyXMini0785.LengthQuantity}
  A nonnegative physical length
\end{definition}
\begin{proof}
  This is the declared model definition of Length Quantity.
\end{proof}
\begin{definition}[Acceleration Quantity]
  \label{def:physics:phyx-mini-0785:phyxminiproblems-problemphyxmini0785-accelerationquantity}
  \lean{PhyXMiniProblems.ProblemPhyXMini0785.AccelerationQuantity}
  A nonnegative linear-acceleration magnitude
\end{definition}
\begin{proof}
  This is the declared model definition of Acceleration Quantity.
\end{proof}
\begin{definition}[Force Quantity]
  \label{def:physics:phyx-mini-0785:phyxminiproblems-problemphyxmini0785-forcequantity}
  \lean{PhyXMiniProblems.ProblemPhyXMini0785.ForceQuantity}
  A nonnegative force magnitude
\end{definition}
\begin{proof}
  This is the declared model definition of Force Quantity.
\end{proof}
\begin{definition}[Moment Of Inertia Quantity]
  \label{def:physics:phyx-mini-0785:phyxminiproblems-problemphyxmini0785-momentofinertiaquantity}
  \lean{PhyXMiniProblems.ProblemPhyXMini0785.MomentOfInertiaQuantity}
  A scalar moment of inertia about the fixed axle, of dimension M L²
\end{definition}
\begin{proof}
  This is the declared model definition of Moment Of Inertia Quantity.
\end{proof}
\begin{definition}[Angular Acceleration Quantity]
  \label{def:physics:phyx-mini-0785:phyxminiproblems-problemphyxmini0785-angularaccelerationquantity}
  \lean{PhyXMiniProblems.ProblemPhyXMini0785.AngularAccelerationQuantity}
  An angular-acceleration magnitude, with radians treated as dimensionless
\end{definition}
\begin{proof}
  This is the declared model definition of Angular Acceleration Quantity.
\end{proof}
\begin{definition}[mass Readout]
  \label{def:physics:phyx-mini-0785:phyxminiproblems-problemphyxmini0785-massreadout}
  \lean{PhyXMiniProblems.ProblemPhyXMini0785.massReadout}
  \uses{def:physics:phyx-mini-0785:phyxminiproblems-problemphyxmini0785-massquantity}
  Read a physical mass in a selected mass unit
\end{definition}
\begin{proof}
  This is the declared model definition of mass Readout.
\end{proof}
\begin{definition}[length Readout]
  \label{def:physics:phyx-mini-0785:phyxminiproblems-problemphyxmini0785-lengthreadout}
  \lean{PhyXMiniProblems.ProblemPhyXMini0785.lengthReadout}
  \uses{def:physics:phyx-mini-0785:phyxminiproblems-problemphyxmini0785-lengthquantity}
  Read a physical length in a selected length unit
\end{definition}
\begin{proof}
  This is the declared model definition of length Readout.
\end{proof}
\begin{definition}[acceleration Readout]
  \label{def:physics:phyx-mini-0785:phyxminiproblems-problemphyxmini0785-accelerationreadout}
  \lean{PhyXMiniProblems.ProblemPhyXMini0785.accelerationReadout}
  \uses{def:physics:phyx-mini-0785:phyxminiproblems-problemphyxmini0785-accelerationquantity}
  Read acceleration in coherent selected length and time units
\end{definition}
\begin{proof}
  This is the declared model definition of acceleration Readout.
\end{proof}
\begin{definition}[force Readout]
  \label{def:physics:phyx-mini-0785:phyxminiproblems-problemphyxmini0785-forcereadout}
  \lean{PhyXMiniProblems.ProblemPhyXMini0785.forceReadout}
  \uses{def:physics:phyx-mini-0785:phyxminiproblems-problemphyxmini0785-forcequantity}
  Read force in coherent selected mass, length, and time units
\end{definition}
\begin{proof}
  This is the declared model definition of force Readout.
\end{proof}
\begin{definition}[moment Of Inertia Readout]
  \label{def:physics:phyx-mini-0785:phyxminiproblems-problemphyxmini0785-momentofinertiareadout}
  \lean{PhyXMiniProblems.ProblemPhyXMini0785.momentOfInertiaReadout}
  \uses{def:physics:phyx-mini-0785:phyxminiproblems-problemphyxmini0785-momentofinertiaquantity}
  Read a moment of inertia in coherent selected mass and length units
\end{definition}
\begin{proof}
  This is the declared model definition of moment Of Inertia Readout.
\end{proof}
\begin{definition}[angular Acceleration Readout]
  \label{def:physics:phyx-mini-0785:phyxminiproblems-problemphyxmini0785-angularaccelerationreadout}
  \lean{PhyXMiniProblems.ProblemPhyXMini0785.angularAccelerationReadout}
  \uses{def:physics:phyx-mini-0785:phyxminiproblems-problemphyxmini0785-angularaccelerationquantity}
  Read angular acceleration in inverse square units of the selected time unit
\end{definition}
\begin{proof}
  This is the declared model definition of angular Acceleration Readout.
\end{proof}
\begin{definition}[mass In Kilograms]
  \label{def:physics:phyx-mini-0785:phyxminiproblems-problemphyxmini0785-massinkilograms}
  \lean{PhyXMiniProblems.ProblemPhyXMini0785.massInKilograms}
  \uses{def:physics:phyx-mini-0785:phyxminiproblems-problemphyxmini0785-massquantity, def:physics:phyx-mini-0785:phyxminiproblems-problemphyxmini0785-massreadout}
  SI kilogram readout of a mass
\end{definition}
\begin{proof}
  This is the declared model definition of mass In Kilograms.
\end{proof}
\begin{definition}[length In Meters]
  \label{def:physics:phyx-mini-0785:phyxminiproblems-problemphyxmini0785-lengthinmeters}
  \lean{PhyXMiniProblems.ProblemPhyXMini0785.lengthInMeters}
  \uses{def:physics:phyx-mini-0785:phyxminiproblems-problemphyxmini0785-lengthquantity, def:physics:phyx-mini-0785:phyxminiproblems-problemphyxmini0785-lengthreadout}
  SI metre readout of a length
\end{definition}
\begin{proof}
  This is the declared model definition of length In Meters.
\end{proof}
\begin{definition}[acceleration In Meters Per Second Squared]
  \label{def:physics:phyx-mini-0785:phyxminiproblems-problemphyxmini0785-accelerationinmeterspersecondsquared}
  \lean{PhyXMiniProblems.ProblemPhyXMini0785.accelerationInMetersPerSecondSquared}
  \uses{def:physics:phyx-mini-0785:phyxminiproblems-problemphyxmini0785-accelerationquantity, def:physics:phyx-mini-0785:phyxminiproblems-problemphyxmini0785-accelerationreadout}
  SI metre-per-second-squared readout of an acceleration
\end{definition}
\begin{proof}
  This is the declared model definition of acceleration In Meters Per Second Squared.
\end{proof}
\begin{definition}[force In Newtons]
  \label{def:physics:phyx-mini-0785:phyxminiproblems-problemphyxmini0785-forceinnewtons}
  \lean{PhyXMiniProblems.ProblemPhyXMini0785.forceInNewtons}
  \uses{def:physics:phyx-mini-0785:phyxminiproblems-problemphyxmini0785-forcequantity, def:physics:phyx-mini-0785:phyxminiproblems-problemphyxmini0785-forcereadout}
  SI newton readout of a force
\end{definition}
\begin{proof}
  This is the declared model definition of force In Newtons.
\end{proof}
\begin{definition}[moment Of Inertia In Kilogram Meters Squared]
  \label{def:physics:phyx-mini-0785:phyxminiproblems-problemphyxmini0785-momentofinertiainkilogrammeterssquared}
  \lean{PhyXMiniProblems.ProblemPhyXMini0785.momentOfInertiaInKilogramMetersSquared}
  \uses{def:physics:phyx-mini-0785:phyxminiproblems-problemphyxmini0785-momentofinertiaquantity, def:physics:phyx-mini-0785:phyxminiproblems-problemphyxmini0785-momentofinertiareadout}
  SI kilogram-metre-squared readout of a moment of inertia
\end{definition}
\begin{proof}
  This is the declared model definition of moment Of Inertia In Kilogram Meters Squared.
\end{proof}
\begin{definition}[radians To Degrees]
  \label{def:physics:phyx-mini-0785:phyxminiproblems-problemphyxmini0785-radianstodegrees}
  \lean{PhyXMiniProblems.ProblemPhyXMini0785.radiansToDegrees}
  Convert a dimensionless radian angle to degrees
\end{definition}
\begin{proof}
  This is the declared model definition of radians To Degrees.
\end{proof}
\begin{definition}[Along Plane Direction]
  \label{def:physics:phyx-mini-0785:phyxminiproblems-problemphyxmini0785-alongplanedirection}
  \lean{PhyXMiniProblems.ProblemPhyXMini0785.AlongPlaneDirection}
  Direction along the inclined plane
\end{definition}
\begin{proof}
  This is the declared model definition of Along Plane Direction.
\end{proof}
\begin{definition}[Surface Contact Regime]
  \label{def:physics:phyx-mini-0785:phyxminiproblems-problemphyxmini0785-surfacecontactregime}
  \lean{PhyXMiniProblems.ProblemPhyXMini0785.SurfaceContactRegime}
  The block's contact regime with the inclined surface
\end{definition}
\begin{proof}
  This is the declared model definition of Surface Contact Regime.
\end{proof}
\begin{definition}[String Model]
  \label{def:physics:phyx-mini-0785:phyxminiproblems-problemphyxmini0785-stringmodel}
  \lean{PhyXMiniProblems.ProblemPhyXMini0785.StringModel}
  Idealization of the connecting string
\end{definition}
\begin{proof}
  This is the declared model definition of String Model.
\end{proof}
\begin{definition}[String Flywheel Contact]
  \label{def:physics:phyx-mini-0785:phyxminiproblems-problemphyxmini0785-stringflywheelcontact}
  \lean{PhyXMiniProblems.ProblemPhyXMini0785.StringFlywheelContact}
  Contact between the wrapped string and the flywheel
\end{definition}
\begin{proof}
  This is the declared model definition of String Flywheel Contact.
\end{proof}
\begin{definition}[Axle Model]
  \label{def:physics:phyx-mini-0785:phyxminiproblems-problemphyxmini0785-axlemodel}
  \lean{PhyXMiniProblems.ProblemPhyXMini0785.AxleModel}
  Constraint and resistance model at the flywheel axle
\end{definition}
\begin{proof}
  This is the declared model definition of Axle Model.
\end{proof}
\begin{definition}[Figure Object]
  \label{def:physics:phyx-mini-0785:phyxminiproblems-problemphyxmini0785-figureobject}
  \lean{PhyXMiniProblems.ProblemPhyXMini0785.FigureObject}
  Physical objects visible in the supplied bitmap
\end{definition}
\begin{proof}
  This is the declared model definition of Figure Object.
\end{proof}
\begin{definition}[Figure Label]
  \label{def:physics:phyx-mini-0785:phyxminiproblems-problemphyxmini0785-figurelabel}
  \lean{PhyXMiniProblems.ProblemPhyXMini0785.FigureLabel}
  Labels printed in the supplied bitmap
\end{definition}
\begin{proof}
  This is the declared model definition of Figure Label.
\end{proof}
\begin{definition}[Supplied Incline Flywheel Figure]
  \label{def:physics:phyx-mini-0785:phyxminiproblems-problemphyxmini0785-suppliedinclineflywheelfigure}
  \lean{PhyXMiniProblems.ProblemPhyXMini0785.SuppliedInclineFlywheelFigure}
  \uses{def:physics:phyx-mini-0785:phyxminiproblems-problemphyxmini0785-figureobject, def:physics:phyx-mini-0785:phyxminiproblems-problemphyxmini0785-figurelabel}
  A structured transcription of the primary figure
\end{definition}
\begin{proof}
  This is the declared model definition of Supplied Incline Flywheel Figure.
\end{proof}
\begin{definition}[Incline Flywheel Setup]
  \label{def:physics:phyx-mini-0785:phyxminiproblems-problemphyxmini0785-inclineflywheelsetup}
  \lean{PhyXMiniProblems.ProblemPhyXMini0785.InclineFlywheelSetup}
  \uses{def:physics:phyx-mini-0785:phyxminiproblems-problemphyxmini0785-massquantity, def:physics:phyx-mini-0785:phyxminiproblems-problemphyxmini0785-lengthquantity, def:physics:phyx-mini-0785:phyxminiproblems-problemphyxmini0785-accelerationquantity, def:physics:phyx-mini-0785:phyxminiproblems-problemphyxmini0785-forcequantity, def:physics:phyx-mini-0785:phyxminiproblems-problemphyxmini0785-momentofinertiaquantity, def:physics:phyx-mini-0785:phyxminiproblems-problemphyxmini0785-angularaccelerationquantity, def:physics:phyx-mini-0785:phyxminiproblems-problemphyxmini0785-alongplanedirection, def:physics:phyx-mini-0785:phyxminiproblems-problemphyxmini0785-surfacecontactregime, def:physics:phyx-mini-0785:phyxminiproblems-problemphyxmini0785-stringmodel, def:physics:phyx-mini-0785:phyxminiproblems-problemphyxmini0785-stringflywheelcontact, def:physics:phyx-mini-0785:phyxminiproblems-problemphyxmini0785-axlemodel, def:physics:phyx-mini-0785:phyxminiproblems-problemphyxmini0785-suppliedinclineflywheelfigure}
  The physical apparatus and its independent response quantities. The flywheel mass is retained even though its separately supplied axial inertia, rather than a shape model, is what enters the rotational equation
\end{definition}
\begin{proof}
  This is the declared model definition of Incline Flywheel Setup.
\end{proof}
\begin{definition}[Matches Problem Readouts]
  \label{def:physics:phyx-mini-0785:phyxminiproblems-problemphyxmini0785-matchesproblemreadouts}
  \lean{PhyXMiniProblems.ProblemPhyXMini0785.MatchesProblemReadouts}
  \uses{def:physics:phyx-mini-0785:phyxminiproblems-problemphyxmini0785-massinkilograms, def:physics:phyx-mini-0785:phyxminiproblems-problemphyxmini0785-lengthinmeters, def:physics:phyx-mini-0785:phyxminiproblems-problemphyxmini0785-momentofinertiainkilogrammeterssquared, def:physics:phyx-mini-0785:phyxminiproblems-problemphyxmini0785-inclineflywheelsetup}
  Numerical and categorical data stated in the exercise prose
\end{definition}
\begin{proof}
  This is the declared model definition of Matches Problem Readouts.
\end{proof}
\begin{definition}[Matches Supplied Figure]
  \label{def:physics:phyx-mini-0785:phyxminiproblems-problemphyxmini0785-matchessuppliedfigure}
  \lean{PhyXMiniProblems.ProblemPhyXMini0785.MatchesSuppliedFigure}
  \uses{def:physics:phyx-mini-0785:phyxminiproblems-problemphyxmini0785-massinkilograms, def:physics:phyx-mini-0785:phyxminiproblems-problemphyxmini0785-inclineflywheelsetup}
  Objects, labels, and incidences read from the primary image
\end{definition}
\begin{proof}
  This is the declared model definition of Matches Supplied Figure.
\end{proof}
\begin{definition}[Angle Rounds To Displayed Tenth]
  \label{def:physics:phyx-mini-0785:phyxminiproblems-problemphyxmini0785-angleroundstodisplayedtenth}
  \lean{PhyXMiniProblems.ProblemPhyXMini0785.AngleRoundsToDisplayedTenth}
  \uses{def:physics:phyx-mini-0785:phyxminiproblems-problemphyxmini0785-radianstodegrees}
  The printed degree label agrees with the physical angle to its displayed tenth
\end{definition}
\begin{proof}
  This is the declared model definition of Angle Rounds To Displayed Tenth.
\end{proof}
\begin{definition}[Uses Textbook Gravity And Angle Calibration]
  \label{def:physics:phyx-mini-0785:phyxminiproblems-problemphyxmini0785-usestextbookgravityandanglecalibration}
  \lean{PhyXMiniProblems.ProblemPhyXMini0785.UsesTextbookGravityAndAngleCalibration}
  \uses{def:physics:phyx-mini-0785:phyxminiproblems-problemphyxmini0785-accelerationinmeterspersecondsquared, def:physics:phyx-mini-0785:phyxminiproblems-problemphyxmini0785-inclineflywheelsetup, def:physics:phyx-mini-0785:phyxminiproblems-problemphyxmini0785-angleroundstodisplayedtenth}
  The textbook numerical calibration implicit in the recorded choice: standard near-Earth gravity and the 3-4-5 trigonometric components whose acute angle rounds to the printed 36.9°. These are measurement/calibration readouts, not assumptions about the requested tension
\end{definition}
\begin{proof}
  This is the declared model definition of Uses Textbook Gravity And Angle Calibration.
\end{proof}
\begin{definition}[Has Physical Parameters]
  \label{def:physics:phyx-mini-0785:phyxminiproblems-problemphyxmini0785-hasphysicalparameters}
  \lean{PhyXMiniProblems.ProblemPhyXMini0785.HasPhysicalParameters}
  \uses{def:physics:phyx-mini-0785:phyxminiproblems-problemphyxmini0785-massreadout, def:physics:phyx-mini-0785:phyxminiproblems-problemphyxmini0785-lengthreadout, def:physics:phyx-mini-0785:phyxminiproblems-problemphyxmini0785-accelerationreadout, def:physics:phyx-mini-0785:phyxminiproblems-problemphyxmini0785-momentofinertiareadout, def:physics:phyx-mini-0785:phyxminiproblems-problemphyxmini0785-inclineflywheelsetup}
  Positivity and nondegeneracy selecting the stated downhill motion
\end{definition}
\begin{proof}
  This is the declared model definition of Has Physical Parameters.
\end{proof}
\begin{definition}[Satisfies Incline Flywheel Dynamics]
  \label{def:physics:phyx-mini-0785:phyxminiproblems-problemphyxmini0785-satisfiesinclineflywheeldynamics}
  \lean{PhyXMiniProblems.ProblemPhyXMini0785.SatisfiesInclineFlywheelDynamics}
  \uses{def:physics:phyx-mini-0785:phyxminiproblems-problemphyxmini0785-massreadout, def:physics:phyx-mini-0785:phyxminiproblems-problemphyxmini0785-lengthreadout, def:physics:phyx-mini-0785:phyxminiproblems-problemphyxmini0785-accelerationreadout, def:physics:phyx-mini-0785:phyxminiproblems-problemphyxmini0785-forcereadout, def:physics:phyx-mini-0785:phyxminiproblems-problemphyxmini0785-momentofinertiareadout, def:physics:phyx-mini-0785:phyxminiproblems-problemphyxmini0785-angularaccelerationreadout, def:physics:phyx-mini-0785:phyxminiproblems-problemphyxmini0785-inclineflywheelsetup}
  The scalar laws use downhill as positive for the block and the corresponding unwinding rotation as positive for the flywheel: normal balance: N = m g cos θ; kinetic friction: fₖ = μₖ N; translation: m g sin θ - fₖ - T = m a; rotation about O: T r = I α; no slip of the taut string: a = r α. None of these laws states a numerical value for T
\end{definition}
\begin{proof}
  This is the declared model definition of Satisfies Incline Flywheel Dynamics.
\end{proof}
\begin{lemma}[string Tension formula]
  \label{lem:physics:phyx-mini-0785:phyxminiproblems-problemphyxmini0785-stringtension-formula}
  \lean{PhyXMiniProblems.ProblemPhyXMini0785.stringTension_formula}
  \uses{def:physics:phyx-mini-0785:phyxminiproblems-problemphyxmini0785-massreadout, def:physics:phyx-mini-0785:phyxminiproblems-problemphyxmini0785-lengthreadout, def:physics:phyx-mini-0785:phyxminiproblems-problemphyxmini0785-accelerationreadout, def:physics:phyx-mini-0785:phyxminiproblems-problemphyxmini0785-forcereadout, def:physics:phyx-mini-0785:phyxminiproblems-problemphyxmini0785-momentofinertiareadout, def:physics:phyx-mini-0785:phyxminiproblems-problemphyxmini0785-inclineflywheelsetup, def:physics:phyx-mini-0785:phyxminiproblems-problemphyxmini0785-hasphysicalparameters, def:physics:phyx-mini-0785:phyxminiproblems-problemphyxmini0785-satisfiesinclineflywheeldynamics}
  Eliminating the normal force, kinetic friction, linear acceleration, and angular acceleration gives the general tension formula. This is a derived relation, not a field or premise of the physical model
\end{lemma}
\begin{proof}
  The conclusion follows from the governing relations and figure readouts named in the statement of string Tension formula.
\end{proof}
\begin{definition}[Answer Choice]
  \label{def:physics:phyx-mini-0785:phyxminiproblems-problemphyxmini0785-answerchoice}
  \lean{PhyXMiniProblems.ProblemPhyXMini0785.AnswerChoice}
  Labels of the four tension choices printed with the exercise
\end{definition}
\begin{proof}
  This is the declared model definition of Answer Choice.
\end{proof}
\begin{definition}[displayed Tension In Newtons]
  \label{def:physics:phyx-mini-0785:phyxminiproblems-problemphyxmini0785-displayedtensioninnewtons}
  \lean{PhyXMiniProblems.ProblemPhyXMini0785.displayedTensionInNewtons}
  \uses{def:physics:phyx-mini-0785:phyxminiproblems-problemphyxmini0785-answerchoice}
  Tension in newtons printed beside each answer label
\end{definition}
\begin{proof}
  This is the declared model definition of displayed Tension In Newtons.
\end{proof}
\begin{definition}[recorded Dataset Answer]
  \label{def:physics:phyx-mini-0785:phyxminiproblems-problemphyxmini0785-recordeddatasetanswer}
  \lean{PhyXMiniProblems.ProblemPhyXMini0785.recordedDatasetAnswer}
  \uses{def:physics:phyx-mini-0785:phyxminiproblems-problemphyxmini0785-answerchoice}
  Dataset answer metadata, not used as a premise of the mechanics
\end{definition}
\begin{proof}
  This is the declared model definition of recorded Dataset Answer.
\end{proof}
\begin{definition}[Is Unique Closest Displayed Tension]
  \label{def:physics:phyx-mini-0785:phyxminiproblems-problemphyxmini0785-isuniqueclosestdisplayedtension}
  \lean{PhyXMiniProblems.ProblemPhyXMini0785.IsUniqueClosestDisplayedTension}
  \uses{def:physics:phyx-mini-0785:phyxminiproblems-problemphyxmini0785-forceinnewtons, def:physics:phyx-mini-0785:phyxminiproblems-problemphyxmini0785-inclineflywheelsetup, def:physics:phyx-mini-0785:phyxminiproblems-problemphyxmini0785-answerchoice, def:physics:phyx-mini-0785:phyxminiproblems-problemphyxmini0785-displayedtensioninnewtons}
  A displayed choice is uniquely closest to the modeled tension
\end{definition}
\begin{proof}
  This is the declared model definition of Is Unique Closest Displayed Tension.
\end{proof}
% --- Archon declaration topology end ---
% --- Archon physics formalization source end ---
